% chapters/requirements/rf_review.tex
% RF-12: Revisión de exámenes
% ============================================================================

\item \textbf{Revisión de exámenes.} \label{rf:review}
La revisión permite al coordinador definir una ventana temporal durante la cual los alumnos
pueden consultar su instancia corregida y solicitar una recorrección de problemas concretos.
Durante la ventana, los alumnos acceden al \acs{pdf}, las anotaciones y la rúbrica (según
los permisos del examen). Tras el cierre de la ventana, los correctores atienden las
solicitudes de revisión y pueden modificar las calificaciones. En la versión~1.0, cada
examen admite una única revisión.

\begin{functional}

    % RF-12.1 ---------------------------------------------------------------
    \item \textbf{Creación de ventana de revisión.} \label{rf:review:create}
    El sistema debe permitir al coordinador crear una ventana de revisión para un examen,
    definiendo la fecha y hora de inicio y fin.

    \textbf{Entrada:} petición POST al recurso de revisiones del examen con la fecha y hora
    de inicio, la fecha y hora de fin y la opción de notificar a los alumnos.

    \textbf{Proceso:} el sistema valida que el examen tiene todas sus instancias en estado
    \texttt{PUBLISHED} y que no existe ya una revisión para este examen (versión~1.0). Se
    crea la revisión con los datos proporcionados. Si se solicita notificación, se encolan
    las notificaciones a los alumnos con la fecha de revisión.

    \textbf{Salida:} respuesta con los datos de la revisión creada y código
    \acs{http}~201.

    % RF-12.2 ---------------------------------------------------------------
    \item \textbf{Apertura automática de acceso.} \label{rf:review:open}
    El sistema debe abrir automáticamente el acceso de los alumnos a sus instancias al
    inicio de la ventana de revisión.

    \textbf{Entrada:} evento temporal: la hora del sistema alcanza la fecha y hora de inicio
    de la ventana.

    \textbf{Proceso:} el sistema transiciona las instancias publicadas del examen al estado
    \texttt{IN\_REVIEW}. A partir de este momento, los alumnos pueden acceder al \acs{pdf},
    las anotaciones, la rúbrica y las calificaciones de su instancia, según los permisos
    configurados en el examen.

    \textbf{Salida:} instancias en estado \texttt{IN\_REVIEW}, accesibles para los alumnos.

    % RF-12.3 ---------------------------------------------------------------
    \item \textbf{Cierre automático de acceso.} \label{rf:review:close}
    El sistema debe cerrar automáticamente el acceso de los alumnos al finalizar la ventana
    de revisión.

    \textbf{Entrada:} evento temporal: la hora del sistema alcanza la fecha y hora de fin de
    la ventana.

    \textbf{Proceso:} el sistema cierra el acceso de lectura de los alumnos a sus
    instancias. Las instancias que no tengan solicitud de revisión pendiente pasan al estado
    \texttt{FINALIZED}. Las instancias con al menos una solicitud de revisión pendiente
    pasan al estado \texttt{PENDING\_REVIEW}.

    \textbf{Salida:} instancias en estado \texttt{FINALIZED} o \texttt{PENDING\_REVIEW}
    según corresponda.

    % RF-12.4 ---------------------------------------------------------------
    \item \textbf{Solicitud de revisión por alumno.} \label{rf:review:request}
    El sistema debe permitir al alumno solicitar la revisión de uno o varios problemas de su
    instancia durante la ventana de revisión.

    \textbf{Entrada:} petición POST al recurso de solicitudes de revisión de la instancia,
    con la lista de problemas a revisar y un mensaje opcional por cada problema.

    \textbf{Proceso:} el sistema valida que la ventana de revisión está abierta y que el
    alumno es el propietario de la instancia. Se crea una solicitud de revisión por cada
    problema indicado, cada una con el mensaje del alumno. Se genera una única petición que
    crea todas las solicitudes simultáneamente.

    \textbf{Salida:} respuesta con las solicitudes creadas y código \acs{http}~201.

    % RF-12.5 ---------------------------------------------------------------
    \item \textbf{Notificación agrupada de solicitudes al corrector.}
    \label{rf:review:notify}
    El sistema debe notificar a cada corrector las solicitudes de revisión que le
    conciernen, agrupadas en una única notificación.

    \textbf{Entrada:} evento interno disparado por el cierre de la ventana de revisión.

    \textbf{Proceso:} el sistema agrupa las solicitudes de revisión por corrector: para cada
    corrector, recopila todas las solicitudes que afectan a problemas que tiene asignados
    (vía \texttt{AssignmentRule}). Se genera una única notificación por corrector con el
    resumen de las solicitudes (número de alumnos, problemas afectados).

    \textbf{Salida:} notificaciones creadas y, si el corrector tiene habilitadas las
    notificaciones por correo, correos enviados con el resumen.

    % RF-12.6 ---------------------------------------------------------------
    \item \textbf{Recorrección por corrector.} \label{rf:review:regrade}
    El sistema debe permitir al corrector atender las solicitudes de revisión, pudiendo
    modificar las calificaciones y anotaciones de las instancias con solicitudes pendientes.

    \textbf{Entrada:} peticiones estándar de modificación de calificaciones y anotaciones
    sobre las instancias en estado \texttt{PENDING\_REVIEW} que el corrector tiene
    asignadas.

    \textbf{Proceso:} el corrector puede consultar los mensajes de las solicitudes, revisar
    el examen, crear nuevas anotaciones, modificar las existentes y ajustar las
    calificaciones. El flujo de calificación es el mismo que en la corrección inicial
    (concurrencia optimista, rúbrica, asignación manual).

    \textbf{Salida:} calificaciones y anotaciones actualizadas según la recorrección.

    % RF-12.7 ---------------------------------------------------------------
    \item \textbf{Cierre de solicitud de revisión.} \label{rf:review:close_request}
    El sistema debe permitir al corrector cerrar una solicitud de revisión.

    \textbf{Entrada:} petición PATCH al recurso de la solicitud de revisión, marcándola
    como resuelta.

    \textbf{Proceso:} el sistema marca la solicitud como cerrada. Si todas las solicitudes
    de la instancia están cerradas, la instancia pasa al estado \texttt{FINALIZED}.

    \textbf{Salida:} respuesta con la solicitud cerrada.

    % RF-12.8 ---------------------------------------------------------------
    \item \textbf{Notificación de resultado de revisión al alumno.}
    \label{rf:review:result}
    El sistema debe notificar al alumno el resultado de la revisión una vez cerradas todas
    sus solicitudes.

    \textbf{Entrada:} evento interno disparado cuando todas las solicitudes de una instancia
    pasan a estado cerrado.

    \textbf{Proceso:} el sistema genera una notificación para el alumno indicando que la
    revisión ha concluido, incluyendo la información de su nota final.

    \textbf{Salida:} notificación \textit{in-app} creada y, si habilitado, correo
    electrónico enviado.

\end{functional}
